\chapter{The Dirichlet Eta Function as an Exact Binary Bridge}
\label{chap:eta}

\section{Alternating Dirichlet series}

The Dirichlet eta function is defined initially for $\Re s>0$ by the convergent alternating series
\begin{equation}
 \eta(s)=\sum_{n=1}^{\infty}\frac{(-1)^{n-1}}{n^s}.
 \label{eq:eta-def}
\end{equation}
For $\Re s>1$, absolute convergence permits an odd-even decomposition:
\begin{align}
 \eta(s)
 &=\sum_{n\ge1}n^{-s}-2\sum_{m\ge1}(2m)^{-s}\notag\\
 &=\zeta(s)-2^{1-s}\zeta(s)\notag\\
 &=(1-2^{1-s})\zeta(s).
 \label{eq:eta-zeta}
\end{align}
Analytic continuation extends this identity to all $s$. The zero of the prefactor at $s=1$ cancels the simple pole of $\zeta$, so $\eta$ is entire.

At $s=1$, the alternating harmonic series gives
\begin{equation}
 \eta(1)=1-\frac12+\frac13-\frac14+\cdots=\ln2.
 \label{eq:eta-one}
\end{equation}
Equations \cref{eq:eta-zeta,eq:eta-one} establish an exact, classical relation among binary alternation, the constant $\ln2$, and the zeta function. This is stronger than a numerical coincidence.

\section{Zero equivalence in the open critical strip}

The binary prefactor vanishes when
\begin{equation}
 2^{1-s}=1.
 \label{eq:prefactor-zero-eq}
\end{equation}
Using $2^{1-s}=\exp((1-s)\ln2)$, its solutions are
\begin{equation}
 s_k=1-\frac{2\pi i k}{\ln2},\qquad k\in\mathbb Z.
 \label{eq:prefactor-zeros}
\end{equation}
For $k\neq0$, these are genuine zeros of $\eta$ lying on $\Re s=1$. At $k=0$, the pole of $\zeta$ cancels the prefactor zero and leaves $\eta(1)=\ln2\neq0$.

Because every prefactor zero has real part $1$, the prefactor is non-zero throughout the open critical strip.

\begin{proposition}[Eta--zeta zero equivalence in $\Sstrip$]
\label{prop:eta-zeta-equivalence}
For $0<\Re s<1$,
\[
 \eta(s)=0\quad\Longleftrightarrow\quad\zeta(s)=0.
\]
The multiplicities agree.
\end{proposition}
\begin{proof}
In $\Sstrip$, the entire factor $1-2^{1-s}$ is non-zero. Multiplication by a non-vanishing analytic function preserves zeros and multiplicities.
\end{proof}

This proposition supplies the most direct exact binary bridge in the initial programme. Every non-trivial zeta zero is a zero of a binary alternating Dirichlet series. The bridge is exact, but it remains many-channel: the alternating series contains infinitely many complex terms. The two-channel cancellation lemma cannot be applied without a canonical reduction.

\begin{figure}[htbp]
 \centering
 \includegraphics[width=0.84\textwidth]{eta_prefactor.png}
 \caption{The eta prefactor has its additional zeros on $\Re s=1$, not on the critical line. In the open strip it is non-vanishing, so eta and zeta zeros coincide there.}
 \label{fig:eta-prefactor}
\end{figure}

\section{Parity channels and a no-go observation}

For $\Re s>1$, define the odd and even channels
\begin{align}
 O(s)&=\sum_{m=1}^{\infty}(2m-1)^{-s}=(1-2^{-s})\zeta(s),\label{eq:odd-channel}\\
 E(s)&=\sum_{m=1}^{\infty}(2m)^{-s}=2^{-s}\zeta(s).
 \label{eq:even-channel}
\end{align}
Then $\eta=O-E$. A candidate identification of an eta zero with equality of these two channels fails after analytic continuation: at a non-trivial zeta zero both $O$ and $E$ vanish because both contain the factor $\zeta(s)$. The equality is then $0=0$, not a normalized cancellation of two non-zero amplitudes.

Away from zeros, their ratio is
\begin{equation}
 \frac{E(s)}{O(s)}=\frac{1}{2^s-1}.
 \label{eq:parity-ratio}
\end{equation}
Equal channel magnitudes would require
\[
 |2^s-1|=1.
\]
Writing $2^s=2^\sigma e^{it\ln2}$, this becomes
\begin{equation}
 2^\sigma=2\cos(t\ln2),
 \label{eq:parity-equal-modulus}
\end{equation}
where the right side must be positive. This is a family of curves, not the line $\sigma=1/2$. Therefore the naive odd-even split does not yield the critical axis.

\section{Alternation and phase $\pi$}

The sign $(-1)^{n-1}$ is $e^{i\pi(n-1)}$. Thus the eta series contains a fixed phase increment $\pi$ between successive parity classes. The same elementary complex phase appears in the separate two-channel cancellation formula, but equality of the phase factor does not identify the underlying mechanisms:
\begin{itemize}
  \item in the two-channel model, two vectors of equal norm cancel;
  \item in eta, infinitely many terms with alternating signs sum to an entire function;
  \item at a zeta zero, global cancellation occurs across the analytically continued function.
\end{itemize}
A canonical coarse-graining would have to be constructed and proved to preserve zerohood and involution covariance.

\section{The eta value $\ln2$ and the half point}

The exact statements form the pair
\[
 \eta(1)=\ln2=\Hbin(1/2),
\]
and
\[
 \eta(s)=0\iff\zeta(s)=0\quad(s\in\Sstrip).
\]
The first equation relates the alternating harmonic series to the numerical value of balanced binary entropy. The second relates eta to the non-trivial zeta zero set. What is not established is
\[
 \eta(s)=0\quad\Longrightarrow\quad\Hbin(\Re s)=\ln2.
\]
That implication is RH-equivalent in this informational formulation and is not assumed.

\section{A precise open state-map problem}

\begin{problem}[Canonical binary reduction of eta]
Construct, from the eta or theta representation and without using the zero set as input, a normalized two-component vector
\[
 V(s)=\begin{pmatrix}V_0(s)\\V_1(s)\end{pmatrix}
\]
and a non-vanishing scalar factor $g(s)$ on $\Sstrip$ such that:
\begin{enumerate}
  \item $\eta(s)=g(s)(V_0(s)+V_1(s))$;
  \item $|V_0(s)|^2=\Re s$ and $|V_1(s)|^2=1-\Re s$ after normalization;
  \item the construction is covariant under $\J$ and analytic at the level required to preserve multiplicities;
  \item no phase or parameter is fitted to known zeros.
\end{enumerate}
\end{problem}
If such a construction existed and satisfied all four requirements, \cref{lem:exact-cancellation} would force all eta zeros in the strip, hence all non-trivial zeta zeros, onto the half-axis. The construction itself remains open.

\status{The eta identity, value at one, prefactor zeros, and zero equivalence in the strip are exact. The normalized binary state reduction is open and is not supplied by the later scalar quotient/kernel construction.}
